% GENERATED FILE. Edit canonical lessons, then run npm run generate:books.
% canonical-source-hash: fnv1a64:3e67c3b30f0726cf
% canonical-lessons: KA-C363-mussanje, KA-C363-mudu, KA-C363-molagu, KA-C363-horahommu, KA-C363-gocarisu

\chapter{Dusk Falls}
\label{ch:ka-dusk-falls}
\hlchaptermodality{363}

\begin{chapteropening}
\textbf{By the end of this chapter:} \emph{I can say dusk, to rise, to resound, to emerge and to become visible.}

Complete the last lesson of chapter 363: I can say dusk, to rise, to resound, to emerge and to become visible.
\end{chapteropening}

\section[\texorpdfstring{\kn{ಮುಸ್ಸಂಜೆ} (mussañje)}{ಮುಸ್ಸಂಜೆ (mussañje)}]{\kn{ಮುಸ್ಸಂಜೆ} (mussañje) — dusk; twilight}

\label{lesson:KA-C363-mussanje}



\begin{quote}
Before the new one: say the Kannada for to demolish, then the Kannada for a statue.
\end{quote}

\subsection*{You'll want to know: \kn{ಮುಸ್ಸಂಜೆ}}
\textbf{\kn{ಮುಸ್ಸಂಜೆ}} — \emph{mussañje} — \textquotedblleft{}dusk; twilight\textquotedblright{}.

\begin{quote}
\kn{ಮುಸ್ಸಂಜೆಯಲ್ಲಿ} \kn{ಹಕ್ಕಿಗಳು} \kn{ಗೂಡಿಗೆ} \kn{ಹಿಂತಿರುಗುತ್ತವೆ}.
\end{quote}

\begin{quote}
\emph{mussañjeyalli hakkigaḷu gūḍige hintiruguttave.}
\end{quote}

\begin{quote}
At dusk the birds go back to their nests.
\end{quote}

\kn{ಮುನ್}, before, and \kn{ಸಂಜೆ}, evening: the time as evening comes on.

\begin{cousinweb}
\textbf{\kn{ಮುಸ್ಸಂಜೆ}} (\emph{mussañje}) is a Kannada compound of native \emph{mun}, \textquotedblleft{}before\textquotedblright{} (DEDR 5020), and \kn{ಸಂಜೆ}, \textquotedblleft{}evening\textquotedblright{}, borrowed from Sanskrit \emph{saṃdhyā}, \textquotedblleft{}juncture, twilight\textquotedblright{} (Monier-Williams).
\end{cousinweb}

\begin{grammarlens}[title={What you've built}]
One more word for telling what happened.
\end{grammarlens}

\subsection*{Your turn}
\begin{itemize}
\raggedright
  \item \emph{Say it:} \emph{mussañje}
  \item \emph{Say it:} \emph{mussañje}, once more
  \item \emph{Say it:} \emph{mussañjeyalli hakkigaḷu gūḍige hintiruguttave.} — At dusk the birds go back to their nests.
  \item \emph{Recall:} say the Kannada for to demolish, then the Kannada for a statue, then say \emph{mussañje} again
\end{itemize}

\begin{tcolorbox}[breakable,colback=teal!4,colframe=teal!35!black,title={Before you move on}]
What does \kn{ಮುಸ್ಸಂಜೆ} mean? (\textquotedblleft{}Dusk; twilight\textquotedblright{}.) Say it once more.
\end{tcolorbox}
\section[\texorpdfstring{\kn{ಮೂಡು} (mūḍu)}{ಮೂಡು (mūḍu)}]{\kn{ಮೂಡು} (mūḍu) — to rise; to appear}

\label{lesson:KA-C363-mudu}



\begin{quote}
Before the new one: say the Kannada for a statue, then the Kannada for dusk.
\end{quote}

\subsection*{You'll want to know: \kn{ಮೂಡು}}
\textbf{\kn{ಮೂಡು}} — \emph{mūḍu} — \textquotedblleft{}to rise; to appear\textquotedblright{}.

\begin{quote}
\kn{ಬೆಟ್ಟದ} \kn{ಹಿಂದೆ} \kn{ಚಂದ್ರ} \kn{ಮೂಡಿತು}.
\end{quote}

\begin{quote}
\emph{beṭṭada hinde candra mūḍitu.}
\end{quote}

\begin{quote}
The moon rose behind the hill.
\end{quote}

Past (it …): \kn{ಮೂಡಿತು}, \emph{mūḍitu}.

The sun, the moon, a doubt or a smile \kn{ಮೂಡು}: \kn{ಮನಸ್ಸಿನಲ್ಲಿ} \kn{ಪ್ರಶ್ನೆ} \kn{ಮೂಡಿತು}.

\begin{cousinweb}
\textbf{\kn{ಮೂಡು}} (\emph{mūḍu}) is native Dravidian: DEDR 5035 lists Kannada \emph{mūḍu}, \textquotedblleft{}to arise, to originate, to become visible\textquotedblright{}, and \emph{mūḍa}, \textquotedblleft{}the east\textquotedblright{}.
\end{cousinweb}

\begin{grammarlens}[title={What you've built}]
One more word for telling what happened.
\end{grammarlens}

\subsection*{Your turn}
\begin{itemize}
\raggedright
  \item \emph{Say it:} \emph{mūḍu}
  \item \emph{Say it:} \emph{mūḍu}, once more
  \item \emph{Say it:} \emph{beṭṭada hinde candra mūḍitu.} — The moon rose behind the hill.
  \item \emph{Recall:} say the Kannada for a statue, then the Kannada for dusk, then say \emph{mūḍu} again
\end{itemize}

\begin{tcolorbox}[breakable,colback=teal!4,colframe=teal!35!black,title={Before you move on}]
What does \kn{ಮೂಡು} mean? (\textquotedblleft{}To rise; to appear\textquotedblright{}.) Say it once more.
\end{tcolorbox}
\section[\texorpdfstring{\kn{ಮೊಳಗು} (moḷagu)}{ಮೊಳಗು (moḷagu)}]{\kn{ಮೊಳಗು} (moḷagu) — to resound; to boom}

\label{lesson:KA-C363-molagu}



\begin{quote}
Before the new one: say the Kannada for dusk, then the Kannada for to rise.
\end{quote}

\subsection*{You'll want to know: \kn{ಮೊಳಗು}}
\textbf{\kn{ಮೊಳಗು}} — \emph{moḷagu} — \textquotedblleft{}to resound; to boom\textquotedblright{}.

\begin{quote}
\kn{ದೇವಸ್ಥಾನದಲ್ಲಿ} \kn{ಗಂಟೆ} \kn{ಮೊಳಗಿತು}.
\end{quote}

\begin{quote}
\emph{dēvasthānadalli gaṇṭe moḷagitu.}
\end{quote}

\begin{quote}
The bell rang out at the temple.
\end{quote}

Past (it …): \kn{ಮೊಳಗಿತು}, \emph{moḷagitu}.

Bells, drums, thunder and slogans \kn{ಮೊಳಗು}.

\begin{cousinweb}
\textbf{\kn{ಮೊಳಗು}} (\emph{moḷagu}) is native Dravidian: DEDR 4989 lists the older Kannada form of this verb, \textquotedblleft{}to sound as musical instruments, to roar, to thunder\textquotedblright{}.
\end{cousinweb}

\begin{grammarlens}[title={What you've built}]
One more word for telling what happened.
\end{grammarlens}

\subsection*{Your turn}
\begin{itemize}
\raggedright
  \item \emph{Say it:} \emph{moḷagu}
  \item \emph{Say it:} \emph{moḷagu}, once more
  \item \emph{Say it:} \emph{dēvasthānadalli gaṇṭe moḷagitu.} — The bell rang out at the temple.
  \item \emph{Recall:} say the Kannada for dusk, then the Kannada for to rise, then say \emph{moḷagu} again
\end{itemize}

\begin{tcolorbox}[breakable,colback=teal!4,colframe=teal!35!black,title={Before you move on}]
What does \kn{ಮೊಳಗು} mean? (\textquotedblleft{}To resound; to boom\textquotedblright{}.) Say it once more.
\end{tcolorbox}
\section[\texorpdfstring{\kn{ಹೊರಹೊಮ್ಮು} (horahommu)}{ಹೊರಹೊಮ್ಮು (horahommu)}]{\kn{ಹೊರಹೊಮ್ಮು} (horahommu) — to emerge; to come out}

\label{lesson:KA-C363-horahommu}



\begin{quote}
Before the new one: say the Kannada for to rise, then the Kannada for to resound.
\end{quote}

\subsection*{You'll want to know: \kn{ಹೊರಹೊಮ್ಮು}}
\textbf{\kn{ಹೊರಹೊಮ್ಮು}} — \emph{horahommu} — \textquotedblleft{}to emerge; to come out\textquotedblright{}.

\begin{quote}
\kn{ಕತ್ತಲೆಯಿಂದ} \kn{ಒಂದು} \kn{ಬೆಳಕು} \kn{ಹೊರಹೊಮ್ಮಿತು}.
\end{quote}

\begin{quote}
\emph{kattaleyinda ondu beḷaku horahommitu.}
\end{quote}

\begin{quote}
A light came out of the darkness.
\end{quote}

Past (it …): \kn{ಹೊರಹೊಮ್ಮಿತು}, \emph{horahommitu}.

Also of results and talents: \kn{ಹೊಸ} \kn{ಆಟಗಾರ} \kn{ಹೊರಹೊಮ್ಮಿದನು}.

\begin{cousinweb}
\textbf{\kn{ಹೊರಹೊಮ್ಮು}} (\emph{horahommu}) is built on native \emph{hora}, \textquotedblleft{}outside\textquotedblright{} (DEDR 4333), and \emph{hommu}, \textquotedblleft{}to rise, to spring up, to come forth\textquotedblright{}, which DEDR 4482 lists only with a query, so its source is not settled.
\end{cousinweb}

\begin{grammarlens}[title={What you've built}]
One more word for telling what happened.
\end{grammarlens}

\subsection*{Your turn}
\begin{itemize}
\raggedright
  \item \emph{Say it:} \emph{horahommu}
  \item \emph{Say it:} \emph{horahommu}, once more
  \item \emph{Say it:} \emph{kattaleyinda ondu beḷaku horahommitu.} — A light came out of the darkness.
  \item \emph{Recall:} say the Kannada for to rise, then the Kannada for to resound, then say \emph{horahommu} again
\end{itemize}

\begin{tcolorbox}[breakable,colback=teal!4,colframe=teal!35!black,title={Before you move on}]
What does \kn{ಹೊರಹೊಮ್ಮು} mean? (\textquotedblleft{}To emerge; to come out\textquotedblright{}.) Say it once more.
\end{tcolorbox}
\section[\texorpdfstring{\kn{ಗೋಚರಿಸು} (gōcarisu)}{ಗೋಚರಿಸು (gōcarisu)}]{\kn{ಗೋಚರಿಸು} (gōcarisu) — to become visible; to show up}

\label{lesson:KA-C363-gocarisu}



\begin{quote}
Before the new one: say the Kannada for to resound, then the Kannada for to emerge.
\end{quote}

\subsection*{You'll want to know: \kn{ಗೋಚರಿಸು}}
\textbf{\kn{ಗೋಚರಿಸು}} — \emph{gōcarisu} — \textquotedblleft{}to become visible; to show up\textquotedblright{}.

\begin{quote}
\kn{ಮಂಜು} \kn{ಕರಗಿದ} \kn{ಮೇಲೆ} \kn{ಬೆಟ್ಟ} \kn{ಗೋಚರಿಸಿತು}.
\end{quote}

\begin{quote}
\emph{mañju karagida mēle beṭṭa gōcarisitu.}
\end{quote}

\begin{quote}
After the mist melted away, the hill came into view.
\end{quote}

Past (it …): \kn{ಗೋಚರಿಸಿತು}, \emph{gōcarisitu}.

A bookish verb; in talk people say \kn{ಕಾಣಿಸು}.

\begin{cousinweb}
\textbf{\kn{ಗೋಚರಿಸು}} (\emph{gōcarisu}) is built on Sanskrit \emph{gocara}, \textquotedblleft{}pasture ground\textquotedblright{}, and so \textquotedblleft{}range, scope, what lies within the senses\textquotedblright{} (Monier-Williams), with the Kannada verb ending -\emph{isu}.
\end{cousinweb}

\begin{grammarlens}[title={What you've built}]
One more word for telling what happened.
\end{grammarlens}

\subsection*{Your turn}
\begin{itemize}
\raggedright
  \item \emph{Say it:} \emph{gōcarisu}
  \item \emph{Say it:} \emph{gōcarisu}, once more
  \item \emph{Say it:} \emph{mañju karagida mēle beṭṭa gōcarisitu.} — After the mist melted away, the hill came into view.
  \item \emph{Recall:} say the Kannada for to resound, then the Kannada for to emerge, then say \emph{gōcarisu} again
\end{itemize}

\begin{tcolorbox}[breakable,colback=teal!4,colframe=teal!35!black,title={Before you move on}]
What does \kn{ಗೋಚರಿಸು} mean? (\textquotedblleft{}To become visible; to show up\textquotedblright{}.) Say it once more.
\end{tcolorbox}
